\chapter{The Constructive Boundary}

The previous chapter wrote a record and then admitted that the record leaned on
promises: a limit the refinements were trusted to approach, a completion marked as
a slot for a later import. This chapter refuses to let the record lean. It replaces
each borrowed analytic existence with a finite artifact the apparatus builds and
holds. Where the earlier chapters said ``a limit exists,'' this chapter produces
the object; where they said ``the residual tends to zero,'' this chapter shows the
residual becomes exactly zero and constructs the tower along which it does. The
constructive boundary of the title is the line between what an apparatus assumes and
what it builds, and the whole chapter is an exercise in moving claims across that
line from the first side to the second.

The chapter opens by dissolving a door it never needed to open. The earlier
convergence was always stated as though it would one day require a theorem from real
analysis --- a Weierstrass density, a completeness, a sup-norm. The first section
shows there is no such door, because the apparatus never instantiated the field that
would have it. Its notion of smallness is order-theoretic: a sequence tends to zero
when it eventually falls under every admissible tolerance. And the convergence the
finite framework actually achieves is stronger than approximation. It is eventual
exactness: once a refinement captures a configuration, the residual there is not
small but zero. ``Weierstrass'' becomes an instantiation of an order fact, not a
cathedral the apparatus must enter.

The second section makes the constructive boundary itself into something the
apparatus can certify, and turns the certification into a method. A construction is
genuinely constructive when the apparatus can run it and read a number off the end,
because a successful run is a certificate that nothing non-constructive was used.
The apparatus runs its integrator, gets a number, and thereby certifies that the
whole path behind the number is built from finite, executable material. It also
isolates the single non-constructive decision the book ever needed --- one classical
choice, sitting at one rung --- so that the rest of the structure can be certified
free of it. Constructed artifacts, produced and run rather than assumed, become the
apparatus's boundary witnesses: where analysis would assert that an object exists,
the apparatus exhibits one and runs it.

The third section proves a fact about the shape of the convergence claim that no
amount of construction can change, and it is a fact of pure counting. The
configurations the apparatus must eventually capture are infinitely many; each
refinement stage is one finite list. So every finite stage misses some
configuration --- more pigeons than holes --- and no single stage can ever support
them all. The quantifier order is therefore forced: not ``there is a stage that
captures every configuration,'' which is provably false for any tower whatever, but
``for every configuration there is a stage that captures it.'' The apparatus does
not get to choose the convenient order; the pigeonhole chooses it.

The fourth section turns the same counting around and extracts the positive result
hiding inside it. The residual does not see a configuration whole; it sees it only
through two cost readings, and those readings take finitely many values. So the
infinite family of configurations collapses, for the residual's purposes, to a
handful of representatives --- three holes --- and solving the residual on the three
representatives solves it on every configuration that shares their tags. The same
counting that forbade a uniform stage hands back a uniform answer up to the cost
quotient. The section also records an honest refusal: one natural action cannot be
solved on all three representatives at once, so no tower can carry it, and for that
action eventual exactness is false rather than merely unproven. The construction
works exactly where the residual already vanishes, and the apparatus says so.

The fifth section produces the two artifacts the chapter has been promising and then
reads the name off the result. It builds a concrete exhausting tower --- an
enumeration of the buildable configurations, graded so that each appears by a stage
the apparatus can compute --- and it builds the convergence certificate as a
constructor rather than a hypothesis, proving the bound from the combinatorics
instead of assuming it. With the tower produced and the certificate proven, the
section reveals what the tower was squeezing all along: the residual the apparatus
has driven to zero is the Fréchet derivative of the action. The derivative exists,
because one linear gauge approximates the action at every path with remainder
exactly zero; it is therefore unique, because any approximating gauge must carry the
same linear part; and the expansion that produced it terminates at order two with no
remainder. The variational principle is not approximately solved along the produced
tower. It is solved.

By the end of the chapter the record of the previous chapter no longer leans. The
limit is a produced tower; the completion's work, where the record needs it, is done
by a constructed artifact; the convergence is eventual exactness proven from
counting; and the quantity at the center of it all has its right name. What the
apparatus could only assume, it now builds.

\section{Eventual Exactness instead of Classical Weierstrass}

The convergence the earlier chapters invoked looks, at first, as if it is waiting
for a theorem from real analysis: a density result, a completeness theorem, a
sup-norm estimate that allows the finite apparatus to borrow a limit. That is the
tempting door. This section closes it. The apparatus does not need to enter a real
line it has not built, and it does not need to pretend that a continuum norm has been
quietly waiting behind the finite record. The convergence needed here has a smaller
and sharper form. It is an order statement inside a scale the apparatus already
carries.

The scale has only four pieces. It has a type of magnitudes. It has a distinguished
zero. It has an order, so one magnitude can be below another. And it has a predicate
that says which magnitudes count as admissible tolerances. That is the whole
structure. There is no completeness, no supremum, no hidden field, and no ambient
metric space. When the apparatus says that a residual tends to zero in such a scale,
it can mean only this: for every admissible tolerance, there is a stage after which
all later residual magnitudes lie below that tolerance,
\[
  \forall\,\tau\ \text{admissible},\ \exists\,N,\ \forall\,n\ge N,\quad
  r_n \le \tau .
\]
This is not a weakened analytic theorem. It is the exact language the finite record
has earned.

The key strengthening is eventual zero. A sequence may approach zero forever without
ever becoming zero; that is the picture suggested by ordinary analytic
approximation. The constructive boundary asks for something stricter. A residual
sequence is eventually zero when there is a stage after which every later residual is
equal to the scale's zero,
\[
  \exists\,N,\ \forall\,n\ge N,\quad r_n = 0 .
\]
If zero lies below every admissible tolerance, then eventual zero immediately gives
tending to zero. After the zero stage, each residual is exactly zero; zero is below
the tolerance; transitivity is enough. No density theorem is being smuggled in. The
argument is order-theoretic from end to end.

This is the chapter's replacement for the old Weierstrass-shaped promise. The finite
record does not claim that residuals become arbitrarily small because an analytic
space has a dense family of approximants. It claims that, once a configuration is
captured by the refinement, the terminal residual for that configuration is exactly
zero. Capture is therefore not approximation. It is arrival. The residual does not
creep toward the target; the apparatus reaches a stage where the target residual is
gone, and all later stages keep it gone.

The natural-number scale is the simple witness that this is not a one-point trick.
Magnitudes can be ordinary natural numbers, zero can be \(0\), the order can be the
usual order, and every natural number can serve as a tolerance. In that scale, the
sample residual that reads \(1\) at the first stage and \(0\) thereafter is already
the whole story in miniature. It is eventually zero from stage \(1\), so it tends to
zero in the scale. The example is small on purpose. It shows that the apparatus can
run a nontrivial scale, a nonconstant residual, and the eventual-zero argument
without asking for a continuum object.

The terminal-residual tower uses the same pattern. Under the explicit exhaustion
hypothesis, each variation's slip is eventually present in the support of the
refinement stage. Once it is present, the carried trace solves the terminal residual
on that variation, so the terminal residual is \(0\). Therefore the sequence of
terminal residual magnitudes over the natural-number scale is eventually zero, and
therefore tends to zero. The old assumed bound is replaced by a finite obligation:
show that the tower eventually supports the variation in question. Where exhaustion
holds, convergence follows because exactness eventually holds.

The hypothesis matters. This section does not yet produce the concrete exhausting
tower. It names the order in which that later construction must pay the debt. Given
a tower that exhausts the relevant configurations, the residual convergence is no
longer a borrowed analytic certificate. It is discharged by the fact that supported
configurations have terminal residual zero. The later sections still have to produce
the tower and show where the exhaustion succeeds. This section fixes the target:
produce support, and the convergence claim becomes eventual exactness.

The scope is therefore honest and narrow. This is not classical real-line
sup-norm density. It is not a claim that every analytic approximation problem has
been solved by counting. It is the order-theoretic convergence appropriate to this
finite apparatus: magnitudes, zero, order, tolerances, and a residual that becomes
zero once refinement captures what it must carry. The square-root and completion
doors remain doors. The apparatus has not crossed them. It has only removed one
false dependency by noticing that the first physical record did not need that door
to make its convergence precise.

The section's anchor is that exactness is humbler and stronger than approximation.
Approximation promises to improve without necessarily arriving. Eventual exactness
arrives and then stops spending language. The apparatus does not chase a limit
through a space it has not built. It states the scale it owns, names the exhaustion
it still owes, and records the consequence: once support is reached, the terminal
residual is not small but zero.

\section{Constructed Artifacts as Boundary Witnesses}

The constructive boundary should not be a matter of the apparatus's good
intentions. It should be a boundary the apparatus can certify in public, by the
same finite discipline it applies to its records. A claim that the apparatus is
constructive is cheap if it is only a label on a paragraph. The claim becomes
serious when the apparatus exhibits the thing, sends it through its carried
operations, and receives a definite reading back. The boundary is then no longer
a mood or a slogan. It is a line drawn by artifacts that have actually been
produced.

The first certificate is execution. A finite construction that can be run on
finite inputs and returns a definite value has already passed a severe test: the
path from input to reading did not require an invisible choice. If such a choice
had been needed, there would be a gap where the apparatus should have carried an
operation. Instead the apparatus puts its construction on the bench, turns the
handle it owns, and reads what comes out. The reading is not merely a numerical
answer. It is evidence about the route that produced the answer. The route has a
carrier, comparisons, residuals, and gauges, and the successful reading says that
these pieces are not borrowed from an existence claim hidden upstream.

The residual example is deliberately small because a boundary witness should be
inspectable. At the first stage the residual reads \(1\). Later, once the
support has been captured, the same residual reads \(0\). Nothing in that story
asks the reader to imagine a completed continuum or a hidden limiting object. The
apparatus shows a finite scale, a residual attached to that scale, and two
readings with different meanings. The first reading says: not yet. The later
reading says: captured. This is the same eventual exactness from the preceding
section, now viewed not as a formal convergence shape but as a produced artifact.
The apparatus owns the residual well enough to evaluate it before and after the
capture point.

The example also fixes the order of explanation. The apparatus does not first
declare that some residual must exist and then search for a reassuring picture.
It begins with the finite rule that produces the residual, applies the rule at a
particular stage, and receives the value \(1\). It applies the same rule later
and receives \(0\). The claim is therefore not suspended over the artifact; it
is read from it. This matters because the surrounding chapter is about the
difference between eventual exactness and mere approximation. The residual
example lets the apparatus say, without analytic ornament, what it means to
arrive: the carried quantity has become zero, and the apparatus can point to the
stage at which the reading changes.

The gauge readings give the same certificate in the language of the first
physical record. On the identity path, the gauge reads \(0\). On a kicked path,
where the apparatus deliberately displaces the comparison, the gauge reads \(2\).
The pair of readings matters because it prevents the certificate from becoming a
vacant success. A gauge that always returned the same value would not tell the
record much. Here the unchanged path and the displaced path separate. The zero
reading is not an inability to see; it is the value appropriate to identity. The
nonzero reading is not a later charge or orientation claim; it is only the
finite sign that the apparatus can distinguish the kicked comparison from the
unchanged one. These two numbers are enough for this section. They certify that
the gauge path is executable and nontrivial, without asking Chapter 11 to spend
claims that belong to later geometry.

So the first boundary witness is not a theorem saying that a witness exists. It
is the witness doing work. The apparatus can point to the sample residual and
say: at this stage it reads \(1\), and at a later stage it reads \(0\). It can
point to the gauge and say: on identity it reads \(0\), and on the kicked
comparison it reads \(2\). The reader is not asked to accept a certificate
detached from its object. The object carries the certificate by being run. This
is the standard of evidence the chapter now imposes on its constructive claims.

That standard also explains why an assumed certificate is still a debt. An
assumption may have the right shape. It may state exactly the convergence or
support condition one wishes to use. It may even look indistinguishable, on the
page, from the conclusion one hopes to earn. But if the certificate has not been
produced by the apparatus, it remains outside the record. It is a promissory
note, not a paid bill. A produced artifact is different. It is not a sentence
asserting that support is eventually reached; it is the finite construction from
which the terminal residual becomes zero. It is not a sentence asserting that a
gauge is constructive; it is the gauge returning \(0\) and \(2\) in the two
tests just described.

This distinction is easy to blur, so the book names it plainly. To assume a
certificate is to move the burden elsewhere. To construct the artifact is to
pay the burden inside the apparatus. The first move may be legitimate as a
temporary scaffold; many earlier passages used scaffolds while the architecture
was still being exposed. But the scaffold is not the finished wall. The
constructive boundary is crossed only when the apparatus replaces a certificate
it had been handed with a certificate it can produce. The residual and gauge
examples show the replacement in miniature. They are not decorative examples.
They are the local grammar of payment.

The same discipline governs the completed tower promised by the preceding
section. A tower equipped with an assumed convergence certificate may be useful
as a temporary model, but it does not yet settle the constructive boundary. It
has a certificate in its pocket, not a certificate earned from its own finite
parts. The apparatus must instead produce the tower, expose the support
condition, and derive the terminal residual from that exposure. Only then can
the convergence claim belong to the object rather than to the surrounding
assumption. The residual \(1\) followed by \(0\) is the small version of that
larger obligation: the artifact must carry the transition from not-yet to
captured by its own finite rule.

The second certificate is a dependency audit. Execution says that a particular
artifact can run. The audit says where the classical boundary sits in the larger
classification ladder. Without the audit, the apparatus would know that some
readings came back, but it would not know whether a non-constructive decision
had entered through a side door of the proofs surrounding those readings. The
audit follows the carried dependencies and separates the rungs. The constructive
rungs clear: the counting rungs, the numeric rung, the comparable rung, and the
gauge reading do not require the classical decision. The one classical decision
appears only at the inferred top rung.

That is the narrowness the chapter needs. There is exactly one named classical
needle in this part of the development, not a fog of classical contamination.
Here the needle keeps the Chapter 1 sense: the exposed yes-or-no edge of a finite
decision. The apparatus does not pretend that the needle is constructive. It
does not dissolve the decision into prose. It also does not let the needle stain
everything below it. The audit names where the decision is used and, just as
importantly, where it is not used. The lower rungs remain constructive artifacts.
The top inferred rung carries the classical edge.

This is why the boundary is sharper than a disclaimer. A disclaimer would say,
somewhere, that a classical ingredient appears. The audit says where. It records
that the constructive base and the gauge reading are clean, while the inferred
top rung is not. The difference matters because Chapter 11 is trying to keep the
apparatus honest without making it timid. If one classical decision appeared
anywhere, and the book treated that as contamination everywhere, then the
finite work below the decision would be erased. If the book ignored the decision,
then the classical edge would be hidden. The audit permits neither mistake. It
keeps the built part built and the borrowed part named.

The audit also prevents a second overcorrection. Once a classical needle has
been named, one might be tempted to treat every later reading as suspect. That
would be a different kind of inaccuracy. The identity gauge reading \(0\), the
kicked gauge reading \(2\), and the residual readings \(1\) and \(0\) do not
inherit the top-rung decision merely because they live in the same chapter. They
belong to the constructive base, and the audit records that fact. The apparatus
therefore does not choose between naive purity and wholesale suspicion. It keeps
an account. Clean rungs stay clean; the inferred rung stays marked.

The two certificates therefore play different roles. Execution is local and
positive: this artifact runs, and here are its readings. Dependency audit is
structural and negative: this path does not draw on the classical decision, and
that top rung does. Together they turn the constructive boundary into a carried
feature of the apparatus. The boundary is not guessed from the style of a proof
and not granted by a phrase such as ``explicit construction.'' It is witnessed by
finite readings and by a ledger of dependencies. The apparatus both produces
the artifact and accounts for what the artifact depends on.

The result is a disciplined form of humility. The apparatus does not say that
all classical reasoning has vanished. It says that, in this finite development,
the constructive material reaches through the residual, the gauge, and the
lower classification rungs, while one inferred top rung remains classical. That
is a much smaller and more useful statement than a global purity claim. It is
also stronger than a vague warning, because it comes with artifacts. The sample
residual reads \(1\) and then \(0\). The identity gauge reads \(0\). The kicked
gauge reads \(2\). The audit leaves one named needle at the top. The record can
now say exactly what it has constructed and exactly what it has merely isolated.

This is the section's metaphysical anchor: to construct is to witness. A constructed
artifact is not an illustration of a witness somewhere else. It is the witness
under the apparatus's own operations. An assumed certificate may guide the
search for such an artifact, but it cannot replace the artifact without leaving
the work unpaid. The boundary between the constructive and the assumed is
therefore not a philosophical preference here. It is a carried boundary: a line
the apparatus can point to, certify by operation, audit by dependency, and
isolate to a single decision when it cannot cross it.

\section{The Forced Quantifier}

The convergence claim of the first section has a quantifier order, and that order
might look like a stylistic choice the apparatus made for convenience. It is not.
Counting forces the order before any later construction begins. The true claim is
per-configuration, or per slip: for each slip there is some stage after which that
slip is supported. The false strengthening reverses the order and asks for one
stage that supports all slips at once. That reversal is not merely too ambitious.
It is incompatible with finite support.

The distinction is small enough to hide in grammar, so the apparatus writes both
forms plainly. The true form says:
\[
  \forall\ \text{slip } s,\ \exists\ \text{stage } N,\ \forall\, n \ge N,\quad
  s \text{ is supported at } n .
\]
The false uniform form says:
\[
  \exists\ \text{stage } N,\ \forall\ \text{slip } s,\ \forall\, n \ge N,\quad
  s \text{ is supported at } n .
\]
The first lets the supporting stage depend on the slip. The second demands a
single finite stage from which the whole infinite family is already captured. The
apparatus is allowed to promise the first only if it later gives the promised
stage slip by slip. It is not allowed to smuggle in the second because the second
is tidier.

This is the place where a harmless-looking swap would change the theorem. If the
stage comes after the slip, the apparatus may inspect the slip and answer with a
stage suited to it. If the stage comes before the slip, the apparatus has to
answer all possible slips without seeing which one is being asked about. The
finite record cannot do that. A stage is not a completed universe of support; it
is one finite support list. The quantifier order records that limitation instead
of hiding it behind a smoother sentence.

Set the obstruction out without drama. Each stage support is a finite list. The
family of slips is unbounded. A finite list can hold only finitely many entries,
so some generated slip always stands outside it. The apparatus makes the escape
explicit by fixing one skeletal move and iterating it. Start from the base slip,
apply the one-step skeleton once, then again, then again. Along that spine, a
depth reading counts how many one-steps have been stacked. Different depths give
different slips, so the family does not fold back on itself. The point is not
that the family is mysterious. The point is that it is generated by a completely
finite recipe and still outruns every finite list.

Now take any proposed stage support \(L\). Since \(L\) is a finite list, the
apparatus can inspect the depths of the slips that \(L\) actually contains and
find a maximum depth. Then it chooses the next iterate: one depth beyond that
maximum. This is the escape witness. It is not selected by an existence fog and
not guessed by an outside observer. It is computed from the list itself:
\[
  \text{escape}(L)
  \quad=\quad
  \text{the one-spine slip at depth } \max(\text{depths in }L)+1 .
\]
Every entry of \(L\) has depth at most the maximum. The escape slip has depth one
larger. Therefore it is not in \(L\). This is the finite-list escape argument in
its useful form: not a slogan about infinity, but a recipe for producing the
missing slip from the finite list that was supposed to contain them all.

The depth maximum is the important finite datum. The apparatus need not survey an
infinite family to refute the list. It only surveys the list. Once the largest
depth appearing there is known, the next depth is enough. This is why the escape
is computable in the ordinary sense of the record: the witness is obtained by a
bounded inspection of the proposed support, followed by one more step along the
spine. The list itself tells the apparatus where to look outside it.

The consequence is immediate. No finite stage supports every slip, so the uniform
statement is false for any tower whatsoever:
\[
  \neg\,\big(\exists\ \text{stage } N,\ \forall\ \text{slip } s,\ 
  s \text{ supported at } N\big).
\]
The only possible form is the per-slip one:
\[
  \forall\ \text{slip } s,\ \exists\ \text{stage } N,\ \forall\, n \ge N,\ 
  s \text{ supported at } n .
\]
The quantifiers must come in that order. The apparatus does not first support
everything and then read off each slip. It must support each slip by its own
eventual stage, and that stage may depend on the slip. A later construction that
claims exhaustion must therefore supply a stage for the given slip, not a magic
stage for all slips at once.

This section is only the negative half. It fixes the logical form and kills the
uniform claim. It does not yet produce the positive tower, and it does not say which
quotients make the later support finite enough to use. Those tasks belong to the
next sections. What this section contributes is the rule those later sections must
obey: the tower may be produced, but its promise must be local in the slip. The
apparatus may answer any slip it is handed, eventually; it may not claim a single
finite stage that answers them all.

That negative result is useful, not merely prohibitive. It tells the later
positive construction what shape its success must have. A produced tower should
not be judged by whether it finds one universal final stage, because no such
stage can exist with finite supports. It should be judged by whether every slip
fed to it receives its own eventual support stage. The obstruction therefore
guards the promise from inflation. It keeps the convergence claim exactly as
strong as the finite apparatus can make true.

The section's anchor is that counting dictates form before construction begins.
The apparatus does not get to choose the convenient order of its convergence
claim. A finite stage is a finite list, and the one-spine family has a depth
beyond any such list. The forced quantifier is not a weakness in the result. It is
the result stated in the only order that the finite record permits.

\section{The Three-Hole Quotient and the Discriminating Refusal}

The same counting that forbade a uniform stage hands back a uniform answer, if
the apparatus asks the right question. The previous section showed that one
finite support list cannot literally contain every slip. That is the negative
half. The positive half begins by noticing that the residual does not inspect a
slip whole. It sees only the cost data that enter the two segments around the
middle node. If many slips are indistinguishable to those two cost readings, then
they are the same as far as the residual is concerned.

Start with that kernel. The residual at the middle node depends on a candidate
slip only through two readings: the cost from the left endpoint to the slip, and
the cost from the slip to the right endpoint. Nothing else of the slip survives
into the residual. So two slips the cost cannot distinguish carry the same
residual:
\[
  \text{cost}(\text{left}, s) = \text{cost}(\text{left}, t)
  \ \text{ and }\
  \text{cost}(s, \text{right}) = \text{cost}(t, \text{right})
  \;\Longrightarrow\;
  \text{residual}(s) = \text{residual}(t).
\]
This is the section's first move. It replaces literal support by quotient
support. The apparatus need not carry every slip separately if the residual
cannot tell some slips apart. It may carry one representative for each residual
class, solve the representative, and transfer the zero across the class.

Now impose the finite tag that the earlier construction already uses. The tag has
three possible values. A cost that factors through that tag reads a slip only by
which of the three values it carries. For residual purposes, the infinite family
therefore collapses into three holes. Choose one representative for tag \(0\),
one representative for tag \(1\), and one representative for tag \(2\). If the
residual vanishes on those three representatives, then it vanishes on every slip
in the matching hole:
\[
  \text{residual}(r_0) = \text{residual}(r_1) = \text{residual}(r_2) = 0
  \;\Longrightarrow\;
  \text{residual}(s) = 0 \ \text{ for every slip } s.
\]
The implication is not magic. Given any slip, look at its tag. Replace it by the
representative with that same tag. Because the cost factors through tag, the slip
and the representative have the same two cost readings around the middle node.
The residual congruence above says their residuals are equal. Since the
representative residual is zero, the slip residual is zero. The three
representatives are not pretending to be all slips. They are enough only because
the residual factors through the tag data.

This is why the section speaks of holes rather than samples. A sample would be a
guess that one observed case stands for many unobserved cases. A hole is a class
defined by the readings the residual actually uses. Once the class has been
defined that way, the representative is not chosen for resemblance or convenience.
It is chosen because every member of the hole has the same residual as the
representative. The transfer is therefore exact, not statistical. The apparatus
does not say that most slips in a hole behave like the representative. It says
that, for this residual and this cost, every slip in the hole has the same
residual value.

This is the positive pigeonhole. The uniform answer that the previous section
proved impossible literally is recovered up to a quotient: not one finite list
containing every slip, but three holes whose representatives carry the residual
answer for every slip in their holes. The apparatus has not defeated the
negative result. It has asked a different, honest question. Literal support
cannot be uniform over the whole family. Residual support can be uniform over the
three tag holes when the cost cannot see finer distinctions.

The transfer theorem should be read with its condition attached. Solving three
representatives solves the quotient only when the residual really factors
through the tag/cost data. If an action asks the residual to see more than the
tag, then the quotient is no longer licensed. If the action does factor through
tag, but its residual refuses to vanish on the three tags themselves, then the
quotient faithfully reports that refusal. The quotient is powerful precisely
because it is narrow. It transfers zero; it does not invent zero.

That condition prevents a common overreading. Three holes are enough for this
residual because the residual sees two segment costs and those costs see only
tags. They are not enough for every possible question one might ask about a
configuration. A later geometric or orientation-sensitive question may split the
same three holes into finer distinctions. This section does not license those
claims. It licenses only the residual transfer that has just been stated. Within
that boundary, the quotient is complete: solve the three representatives and the
whole quotient is solved.

The discriminating action shows the limit. It assigns no cost between equal tags
and unit cost between unequal tags. Around a middle node, the residual asks how
the left-to-middle and middle-to-right comparisons change when the middle tag is
varied. Across the three possible middle tags, that pattern is nonconstant. In
one arrangement the values have the shape \(\{0,2,2\}\); in another they have the
shape \(\{1,1,2\}\). The details are finite arithmetic over three tags, but the
consequence is structural: at least one tag representative has nonzero residual.
There is no way to make the residual vanish on all three at once:
\[
  \text{no covering trace makes the discriminating residual vanish on all three tags.}
\]
This is the same arithmetic that produced the kicked reading \(2\): two
boundary-crossing edges. Here it is not promoted into a later sign, charge, or
orientation claim. It is only the finite obstruction the residual sees. Because
the discriminating residual is nonzero somewhere among the three
representatives, no trace covering those representatives can solve it. And
because the discriminating cost itself factors through tag, the refusal is not
limited to those literal representatives. Any knot set that covers all three
tags would transfer a forced zero back to the three representatives, contradicting
the obstruction. Thus there is no covering trace and no tag-covering trace for
this action.

That refusal is not a failure of the quotient method. It is the quotient method
working honestly. The method says: if the residual factors through the three
holes, then solving the holes solves the family. The discriminating action says:
the residual over the three holes is not all zero. Therefore the method refuses
to certify it. The apparatus is not stuck because a proof is missing; it is
informed that this action is not stationary under the residual test. Eventual
exactness is blocked for the right reason.

The no-covering conclusion has two layers. First, a trace that literally carries
the three representatives would force the terminal residual to vanish on those
representatives, contradicting the nonzero tag pattern. Second, even a trace that
uses different knots cannot escape if those knots cover the same three tags. Tag
congruence would carry the forced zero from the chosen knot back to the
representative of its tag. The obstruction therefore survives replacement of the
representatives. It is a tag-level obstruction, not an artifact of a bad choice
of three named slips.

What can pass through the quotient is an action whose residual is structurally
zero. The clean repair is telescoping. Its cost is read as one tag plus the
complement of the other:
\[
  \text{cost}(x,y)=\mathrm{tag}(x)+(2-\mathrm{tag}(y)).
\]
Around any middle node the two segment costs telescope:
\[
  \text{cost}(\text{left}, t) + \text{cost}(t, \text{right})
  = \mathrm{tag}(\text{left}) + 4 - \mathrm{tag}(\text{right}),
\]
which does not depend on the middle \(t\) at all. So the difference under any
variation is identically zero. The action is nontrivial as a cost, but its
residual is structurally zero. It is flat at every middle node.

There is a small subtlety worth keeping visible. The telescoping action is not
accepted because it behaves like an ordinary path admission rule. Its composed
cost differs from the direct endpoint cost by a constant offset. That would be
the wrong door for an earlier admission test. The present covering route asks a
different question: whether the residual vanishes on the tag representatives and
therefore on the quotient. For the telescoping action it does. The residual is
zero for every path and every variation directly, and the three-hole transfer
also certifies the same fact by the quotient route. This is the constructive
repair: finite tag data, structurally zero residual, flatness at every middle
node, and an accepted covering route for the reason this chapter actually needs.

This also explains why the repair is not a trick that erases the discriminating
refusal. The telescoping action changes the action, not the verdict about the
discriminating one. It arranges the two segment costs so the middle tag cancels
out of the residual. The discriminating action arranged the costs so changing
the middle tag could be detected. The quotient records both facts. One action
cannot pass because its three tag residuals are not all zero. The other passes
because its residual is zero before any representative transfer is needed, and
the transfer then confirms the same result across the quotient.

This pairing is the section's lesson, and it keeps the constructive program
honest. The quotient is powerful but narrow. It certifies eventual exactness
exactly where the residual is zero on the covered representatives, and not one
step further. An action with genuinely nonzero residual is not a gap in the
argument; it is an action the finite gauge equation declares unsolved. The
apparatus that could drive every residual to zero would be an apparatus whose
equation said nothing. The refusal is part of the content.

The scope also stays local. This section supplies the quotient and the
obstruction. It does not yet construct the concrete exhausting tower, and it does
not yet claim Frechet uniqueness. Those belong to the next section. What has been
earned here is the middle result: when cost collapses slips into three residual
holes, three solved representatives transfer zero to the whole quotient; when the
three holes themselves cannot be solved, no quotient language may pretend
otherwise; and when a telescoping action makes the residual structurally zero,
the covering route closes for the right reason.

The section's anchor is that a quotient is honest only when it records both what
it collapses and what it refuses to collapse. The cost sees three holes, and the
apparatus may solve three representatives in place of an infinite family only
because residual equality follows from cost equality. But the same finiteness
also exposes actions that cannot be solved on the three holes at all. The
constructive boundary is not only a line the apparatus works up to; it is also a
line the apparatus can prove it must respect, for actions whose residual will not
vanish.

\section{The Produced Tower and Frechet Uniqueness}

The chapter has promised two things and delayed both until the counting
apparatus could honestly support them. First, it promised a concrete tower that
really exhausts the configurations it claims to exhaust. Second, it promised a
convergence certificate that is produced by the tower rather than smuggled into
the hypothesis. The earlier sections prepared the ground: eventual exactness
replaced classical density, constructed boundary witnesses replaced assumed
certificates, the quantifier order blocked a single universal stage, and the
three-hole quotient separated a passing action from an obstructed one. This
section closes the numbered argument. It produces the tower, produces the
certificate, and then reads the name of the residual it has driven to zero.

The tower does not enumerate the full universe of possible configurations. That
would be the wrong promise. The full universe contains more than the finite
apparatus can list by a natural-number clock. Some of its distinctions are raw
logical or type-level commitments, not generated pieces in a countable
construction. A countable tower cannot cover those by pretending they are
ordinary entries in a finite list. The honest enumerable family is narrower: the
rational skeleton. It begins with a base configuration and closes under the
constructible events over a fixed set of atoms. Every member of this skeleton
has a construction history, and that history supplies the grading.

The grading is depth. A base entry has depth \(0\). A one-step extension has
depth one more than the entry it extends. A two-branch entry has depth one more
than the deeper of its two branches. This rule matters because it makes the
finite slices exact. A stage is not merely a bag of entries with approximately
the right size. Stage \(n\) contains all skeleton entries whose construction
depth has become visible by \(n\), and the tower is cumulative: once an entry
appears, later stages keep it. Thus a skeleton entry \(c\) of depth \(d\) comes
with its own stage, computed from the entry itself:
\[
  \text{depth}(c)=d
  \;\Longrightarrow\;
  c \in \text{stage}(d)
  \;\Longrightarrow\;
  c \in \text{stage}(n)\quad(d\leq n).
\]
This is stronger than a vague refinement relation. It is genuine
entries-monotonicity. Membership persists, so the eventual claim has something
stable to talk about.

The tower therefore exhausts exactly what it says it exhausts. Given any
skeleton slip, compute its depth. That number is the first stage from which the
slip is supported forever. The tower does not need to guess a witness and it
does not ask an outside completion principle to supply one. It computes the
witness from the slip. This is the positive counterpart to the forced
quantifier of the third section:
\[
  \forall s\ \text{in the skeleton}\;\exists N\;\forall n\geq N,\quad
  s\in\text{stage}(n).
\]
The order of quantifiers is still the honest one. The stage \(N\) depends on
the slip \(s\). The tower does not collapse this into a single \(N\) that works
for all slips.

Indeed, the produced tower also preserves the negative result. It supports each
skeleton slip eventually, but no one stage supports all skeleton slips. Given a
finite stage, the one-spine escape from the forced-quantifier section still
applies: choose a generated slip one depth beyond everything in the finite
list. The produced tower has not cheated the pigeonhole. It has supplied the
right per-slip stage while leaving the forbidden uniform stage forbidden. This
is why the construction is trustworthy. It solves the problem the finite
apparatus can solve, and it leaves untouched the problem the counting argument
proved unsolvable.

The same tower also remembers the obstruction from the previous section. Early
finite stages already cover the three tag representatives, and from that point
forward every later stage covers them as well. For the discriminating action,
covering those tags is exactly what cannot be made stationary. The obstruction
therefore survives inside the produced tower. The tower gives the accepted
action a route; it does not rescue the refused action. A nonzero residual is
not a missing certificate waiting to be discovered. It is a genuine failure of
eventual exactness for that action.

The contrast is worth making explicit before the certificate enters. The
rational skeleton is not a retreat from the book's ambition; it is the first
place where the ambition becomes holdable. The full configuration universe
contains more distinctions than a finite tower can index. The skeleton records
the configurations generated by the apparatus's own finite grammar. On that
family, every membership claim has a finite reason: either the entry is the base
entry, or it arises from one of the constructors applied to earlier entries.
Depth is therefore not a decorative rank. It is the receipt of construction. To
know the depth of a skeleton slip is to know how far the tower must climb before
that slip becomes visible.

This also explains why the tower uses cumulative stages rather than disjoint
slices. A disjoint slice would say where an entry first appears, but it would
not by itself say that the entry remains available when later stages perform
their finite work. Eventual support requires persistence. Once the slip has
entered, every later finite stage may use it. That persistence is the practical
form of the word ``eventual'' in this chapter. Eventual exactness is not a
momentary hit at one stage; it is stable support from some stage onward.

Now produce the certificate. An earlier completed tower could carry a field
saying that its bound tends to zero. But a carried assertion is not yet a
produced certificate. If the conclusion follows because the structure was handed
the very convergence it needs, nothing has been discharged. The honest
certificate must construct the bound and prove its behavior from the finite
stage data. It must show where the bound comes from, why it dominates the
residual, and why it tends to zero.

For the telescoping action, the construction is especially clean. At each finite
stage, the trace consumes its entire knot list by identity steps. No entry is
left unprocessed inside that finite stage. The terminal residual is the residual
itself. The bound is chosen to be that same measured residual. Then the
telescoping calculation from the previous section does the decisive work: the
middle term cancels, so the residual is structurally zero at every middle node.
Consequently, at every stage,
\[
  \text{bound}(n,s)=\text{measured residual}(n,s)=0.
\]
The squeeze is not asymptotic theater. The residual is already zero, the bound
equals it, and the constant zero sequence tends to zero in the order-theoretic
scale of the first section.

This is the difference between an assumed certificate and a produced one. An
assumed certificate says, ``suppose the bound tends to zero,'' and then reads
off convergence. The produced certificate says: here is the finite tower, here
is the trace at each stage, here is the terminal residual, here is the bound,
and here is the proof that the bound is zero at every stage. The statement
\[
  \text{bound tends to zero}
\]
is no longer a premise. It is a result of the construction. Eventual exactness
holds in its strongest form: not because the sequence slowly approaches zero,
but because the sequence is constant at zero.

The construction also clarifies the role of the exhaustive tower. For the flat
telescoping action, the residual vanishes before any delicate limiting argument
is needed. The exhaustion route is still important because it proves the
finite-stage apparatus has the right coverage story: each skeleton slip enters
at its computed depth and stays there. But the zero reading is not produced by
waiting long enough for a nonzero quantity to decay. The zero reading is
structural. The tower carries that structural zero through a produced
finite-stage certificate.

There is no hidden analytic compactness in this move. The scale here is the
natural-number stage scale, and convergence means the order-theoretic eventual
zero condition already introduced at the start of the chapter. The produced
certificate proves that condition directly: choose stage \(0\), and every later
stage reads zero. The phrase ``tends to zero'' is therefore not borrowing metric
intuition. It names a finite order fact about a sequence of readings. The
apparatus can inspect every part of the claim: the stage, the terminal
residual, the bound, and the witness that the bound is already zero.

For the discriminating action, the same language gives the opposite verdict.
Its residual is nonzero on at least one tag representative, and any
tag-covering stage would have to face that representative. No tower certificate
can turn that residual into zero without contradicting the obstruction. The
right conclusion is not that the discriminating action has not yet been
approximated well enough. The right conclusion is that eventual exactness is
false for it. The apparatus now has both sides: a produced flat case and a
proved obstruction.

With the tower and certificate in hand, the section can name the quantity the
book has been manipulating. The residual squeezed by the tower is the first
variation of the action. In this finite apparatus, that first variation is the
Frechet derivative. The identification is not an analogy added after the fact.
The same expression that has appeared as the spline residual, the terminal
residual, and the stage reading is the derivative reading:
\[
  \text{residual}
  =
  \text{first variation}
  =
  \text{Frechet derivative}.
\]
The tower was never chasing an auxiliary error term. It was driving the
derivative reading of the action to zero.

Existence of the derivative is also finite and explicit. From an action, the
apparatus constructs a gauge whose first-order part matches the action's composed
difference at every path and every variation, with remainder exactly zero:
\[
  \exists G\;\forall p,v,\quad
  \text{action difference at }(p,v)
  =
  G_{\mathrm{lin}}(p,v)
  \quad\text{and}\quad
  \text{remainder}(p,v)=0.
\]
That is the finite existence theorem. It does not require a normed vector space
of variations, and it does not require a limiting argument over small
increments. The remainder is zero by construction, so the approximating gauge
exists in the strongest possible sense available here.

The phrase ``at every path'' should not be read as a continuum statement. A path
in this apparatus is a finite object with named origin, middle, and target data.
The derivative exists because the action supplies a finite composed difference
against each finite variation, and the canonical gauge records that difference
with no leftover term. The construction is global over the finite path grammar,
not over an analytic manifold. That is why the theorem can be direct without
becoming grandiose.

Uniqueness follows by exact determination. Suppose two gauges approximate the
same action at the same path. Since both match the same composed difference,
their first-order readings at that path must agree on every variation:
\[
  G_{\mathrm{lin}}(p,v)=H_{\mathrm{lin}}(p,v)
  \quad\text{for every }v.
\]
This is the finite version of Frechet uniqueness. In the classical setting,
linearity and a small remainder isolate the derivative. Here the remainder is
identically zero, so the isolation is sharper: two candidates that exactly
match the same difference have no room to disagree. Existence and uniqueness
are not separate mysteries. The construction gives a derivative, and exact
matching pins it down.

This uniqueness is local in the right way. It says that at a fixed path, every
approximating gauge has the same first-order part. It does not say that two
different actions become the same, or that two raw configurations must be
identified. The uniqueness belongs to the derivative reading once the action,
path, and matching condition have been fixed. That boundary keeps the theorem
from turning into an overclaim about the surrounding configuration world.

The same care is needed for the cubic expansion. The apparatus has several
slots: the order-one linearization, the left cubic slot, the right cubic slot,
and the second variation. Each slot has its own matching condition. The
order-one part is pinned by matching the composed difference. The left slot is
pinned by the one-coordinate left difference. The right slot is pinned by the
one-coordinate right difference. The second variation is pinned once the
canonical first-order terms have been fixed. The warning is important: the bare
three-term sum does not determine its terms. A constant could be shifted between
the left and right first-order slots while preserving the total. Slot-by-slot
uniqueness depends on the individual matching conditions, not on the sum alone.

With those matching conditions in place, the discrete Taylor statement
terminates at order two:
\[
  D(v)
  =
  r_{\mathrm L}(v_{\mathrm L})
  +
  r_{\mathrm R}(v_{\mathrm R})
  +
  \delta^2(v),
  \qquad\text{with no remainder.}
\]
This is not an approximation plus an error estimate. The apparatus reaches an
exact finite decomposition. The expansion stops because the finite grammar has
no further undischarged remainder slot in this calculation.

The two first-order slots deserve their separate names. The left slot measures
the one-coordinate change on the left side of the middle node; the right slot
measures the corresponding one-coordinate change on the right side. The second
variation records the coupled remainder after those canonical first-order
contributions have been fixed. If one forgets those matching conditions and
looks only at the displayed sum, the decomposition can be misread as arbitrary.
The apparatus avoids that mistake by assigning each term its own finite test.
The decomposition is exact because each term has been separately identified.

The word ``linear'' must stay honest. There is no ambient vector space of
variations in this argument, so the apparatus does not claim classical
linearity of a map between normed spaces. The role usually played by linearity
is played here by uniqueness under the matching condition. Likewise, the
classical small-\(o\) remainder condition is not imported as analytic structure.
The finite remainder is identically zero, so it satisfies the intended
smallness demand without a norm. The genuine continuum derivative, with
ordinary vector-space linearity and limiting smallness, remains behind the
square-root door of the previous chapter.

That door matters. The finite Frechet derivative here is not a covert claim that
the classical continuum derivative has already been recovered. It is the finite
object occupying the derivative role inside the apparatus: first variation,
exactly matched gauge, zero remainder, and uniqueness under matching. The
continuum question asks for additional structure, and the book has deliberately
kept that additional structure on the other side of the square-root passage.
The present section closes the finite calculus, not the analytic one.

The sequence-level finish is now precise. For the telescoping action, the
measured Frechet derivative along the produced tower tends to zero because it
is zero at every stage. More generally, whenever an action is flat at a path,
the measured derivative reading is zero and the resulting constant sequence
tends to zero. Conversely, if a residual reading is genuinely nonzero, then the
constant sequence of its measured magnitude cannot tend to zero in this scale.
The apparatus therefore records both directions:
\[
  \text{flat derivative readings tend to zero,}
  \qquad
  \text{genuinely nonzero readings do not.}
\]
The variational principle is not merely witnessed in a showcase example. It is
stated as a finite sequence principle with an obstruction clause.

The numbered chapter closes here because the promised finite work has now been
done. The rational skeleton gives a depth-graded cumulative tower. The forced
quantifier survives: every skeleton slip has its own eventual stage, and no
single stage supports them all. The telescoping action receives a produced
certificate whose bound is the residual itself and whose residual is zero at
every stage. The discriminating action remains refused for exactly the same
finite reason it was refused before. The residual is named as the Frechet
derivative; the derivative exists by construction, is unique by exact matching,
and terminates at order two in the finite expansion.

The metaphysical point is that the apparatus produces what it is allowed to
claim. It does not appeal to a hidden continuum, and it does not replace a
missing argument by a field labeled convergence. It constructs a finite tower,
constructs a finite certificate, and names the derivative only after the residual
has been made inspectable at every stage. The finite variational principle is
therefore not a limit the apparatus hopes to approach. It is a constructed
boundary witness: a tower, a zero certificate, and a unique derivative reading
held together by finite count.

\section*{Bridge: A Derivative Told in One Coordinate}

The chapter ends with a clean achievement. It gave the finite tower that the
preceding obstruction demanded, kept the depth grades honest, and produced the
zero certificate that lets the telescoping action count as flat at every
inspected stage. It did not ask the reader to trust an unexamined passage to a
continuum. It built the rational skeletons, accumulated the finite stages,
watched the residual fall to its terminal value, and then named the first
variation as a finite Frechet derivative only after the residual had become a
visible witness. The derivative exists because the tower produces it. It is
unique because exact matching leaves no room for a second reading. Its finite
Taylor expansion terminates because the apparatus has only the slots it has
counted. That is real success, and the chapter should receive it as such.

The success has a shape, though, and the shape matters. This derivative tells a
one-coordinate story. It reads what happens when the action changes along a
single line of variation, and it reads that line exactly. Nothing in the
chapter weakens that reading. The point is narrower and more dangerous: exact
does not mean exhaustive. A key can fit the one lock it was cut for and still
tell us nothing about the neighboring door. The finite derivative of this
chapter fits its coordinate. It measures the first variation in the direction
that the instrument has chosen. It does not thereby measure every quantity that
can appear when two sides of the middle change together. The limit lies in the
question, not in the answer. A perfect answer to a single-coordinate question
still inherits the exclusions that made that question single-coordinate in the
first place.

The earlier coupled difference already warned us about this narrowing. That
difference did not divide into a single first-order story plus a decorative
remainder. It carried a left first-order slot, a right first-order slot, and a
mixed order-two coupling. Those three pieces do not play the same role. The
left slot says what the left side contributes when it changes as itself. The
right slot says the corresponding thing for the right side. The mixed coupling
does something else: it records the residue that belongs to their joint
disturbance. It does not sit inside either first-order slot waiting for a more
patient reader. It appears at the place where the two first-order stories meet
and fail to exhaust the difference between them.

That is why the one-coordinate derivative can be exact and still miss the
coupling. If the left side alone moves, the right side supplies no partner for
the cross term. If the right side alone moves, the same silence returns from
the other side. Each one-coordinate reading can report its own first-order
slot, and each report can be correct. But the mixed term has no voice until the
variation moves as a pair. A single coordinate does not merely overlook it by
accident. The single coordinate gives the mixed term no place to speak. The
instrument asks a question whose grammar has only one changing side, and the
cross coupling answers only questions whose grammar has two.

This also clarifies the role of the weak form that has been so useful. The
weak form lets the chapter isolate a derivative by projecting away what would
otherwise contaminate a one-coordinate reading. That discipline made the zero
certificate legible. It allowed the tower to certify flatness without smearing
the first variation into the coupling. But projection also has a cost. What it
removes does not stop existing just because it no longer appears in the
one-coordinate account. The weak form has helped the apparatus tell the clean
story; the next chapter must ask what the clean story left outside its frame.

The natural name for the next question is measurement-minus-story. The story
says: here is the left first-order reading, here is the right first-order
reading, and here is the finite derivative that each one-coordinate instrument
can report. Measurement asks whether those reports exhaust what the paired
change leaves behind. If they do not, then the residue no longer looks like an
error in the derivative. It looks like the thing that the derivative was never
designed to read. The apparatus must therefore change the question rather than
weaken the answer it just earned.

Chapter 12 begins at that change of question. It does not need to revoke the
finite Frechet derivative, and it should not. The derivative stands: exact,
unique, produced by a finite tower, and terminated by a finite expansion. The
next chapter asks what appears when the variation receives two coordinates at
once. It returns to the coupled difference, lets the two sides move together,
and watches the mixed place that a one-coordinate reading had forced into
silence. There the signed residue becomes the problem to identify, not a
conclusion already in hand. The bridge from this chapter to the next therefore
turns on a precise unease: if the one-coordinate derivative has succeeded on
its own terms, and if those terms still cannot hear the mixed coupling, then
what kind of instrument can measure the residue that belongs only to the pair?

\section*{Coda: The Locksmith's Master Key}

A locksmith at a bench, cutting keys for a long corridor of doors, performs the
constructive discipline of this chapter with brass, pins, and a depth gauge. The
trade is worth watching because it gives the chapter a body. The blade on the
bench has no opinion about convergence, yet the lock judges it with the same
severity that the chapter has been trying to earn. The cut reaches the required
depth or it does not. The pin stack rises to the shear line or it does not. The
cylinder turns or it holds. There is no charitable middle in which a nearly
right key opens a door nearly enough.

A key is therefore not approximated in the sense that matters to the lock. It is
cut to exact depths. The locksmith may approach the depth slowly, may file a
little and check the gauge, may polish the cut before testing the blade; but the
door does not reward the approach. A key ninety percent of the way to the right
depth does not open the lock ninety percent. It fails with the same finality as
a blank. The shear line is the lock's event of exactness. Until the pin stacks
meet it, the mechanism refuses. Once they meet it, the cylinder turns. The
chapter's eventual exactness has this same plain shape: arrival matters, and
approach matters only as the finite work by which arrival is reached.

The lock also reads through a quotient. It does not care about the manufacturer's
stamp, the color of the bow, the scratches left by a pocket, or the story of
who carried the key. It sees the depths. More precisely, it sees the heights to
which those depths lift a finite list of pin stacks. Two blades with the same
depth pattern are the same key to that cylinder, however different they look in
the hand. The whole visible variety of blanks collapses at the plug to a finite
list of readings. The cylinder is a quotient device: ornament vanishes, brand
vanishes, handling marks vanish, and the pin-depth pattern remains.

That is the three-hole lesson in metal. The lock reads the slots it has, not
every distinction the world can name. If a scratch does not change a pin height,
the lock ignores it. If two keys differ in shape outside the read cuts, the
lock treats them as one. If three visible keys lift the same finite pin pattern,
the cylinder does not know that three visible keys arrived. It knows the
pattern. The locksmith relies on this every day. To make a working key, he does
not reproduce the whole visible history of the original. He reproduces the
finite pattern the lock can test.

The forced quantifier appears as soon as the locksmith opens the drawers. A shop
may keep trays of blanks and racks of cut keys, but no single finite tray
already contains every possible key for every possible cylinder. Each lock
pattern asks for its own work order. One door may require three cuts at shallow
depths. Another may require five cuts, with the fourth pin riding high. A third
may need a blank from a different rack before any cut can even begin. For each
lock that belongs to the day's work, the locksmith can find the blank, set the
depth gauge, cut the blade, and place the finished key in the box. But the box
does not become a universal archive. It grows by finite jobs, and each job has
its own stage.

Once a pattern enters the box, it remains available. The locksmith does not
unmake yesterday's key in order to cut today's. The rack is cumulative: a key
cut for one cylinder sits in its labeled place while the next work order
arrives. This is the depth-graded tower in shop language. At shallow depths,
the locksmith has only the simple patterns. Later, the more detailed patterns
enter. Each pattern appears at a finite stage determined by its own demand, and
thereafter it can be taken down and tested again. The promise is not that one
morning's tray held every key. The promise is that each finite pattern, when
its turn arrives, receives a cut blade that stays in the record.

The order on the rack matters as much as the rack itself. A careless shop could
pile keys in a drawer and claim that everything was somewhere inside. The
locksmith's real discipline is sharper. Each hook has a label. Each label names
the finite pattern the hook supports. When a harder cylinder arrives, the
locksmith does not disturb the easier ones; he adds another labeled hook at the
right depth of the rack. If a client returns with an earlier lock, the earlier
key is still there. If a later client brings a more delicate cylinder, the shop
extends the rack rather than pretending the old tray already solved the new
case. That is why the tower is not a heap. It has stages, persistence, and a
visible rule for placing the next artifact where it can be found again.

The master key sharpens the issue. A client may want one blade that opens a
whole group of locks. The wrong way to satisfy that request is to assert that
such a blade must exist and hand over a blank with confidence. The locksmith
has to inspect the cylinders. He has to read the pinning charts, compare the
positions, and find a pattern that each lock in the group will accept. If the
pattern exists, it does not exist as a wish. It exists as brass. The locksmith
cuts it, tests it, and watches each cylinder turn. The master key is the
certificate. Its success lies in the doors it opens, not in a note saying that
some key should open them.

That produced key is the coda's analogue of the produced certificate. The
chapter has refused to let an assumed certificate do the work. It has asked for
the artifact: the finite tower, the zero reading, the trace whose bound is the
residual itself, and the residual that vanishes in the case where the action
truly rides. In the shop, the same distinction is tactile. A promise that a
master key can be cut leaves every door shut. A key that has been cut and tested
turns the cylinders. The difference between promise and proof is the difference
between a blank and a blade whose cuts have met the pins.

The locksmith also knows the dignity of refusal. Some groups of locks do not
admit a master key. If one cylinder demands a high lift at a pin position while
another demands a low lift there, no single cut can satisfy both. The conflict
does not call for more optimism or a finer file. It calls for an honest verdict:
this group cannot share one blade. The refusal is not failure by exhaustion. It
is a successful reading of incompatibility. The pinning charts show the clash,
and the locksmith can certify the refusal before he ruins a blank pretending
otherwise.

That refusal is as constructive as the successful master key. It tells the
client something exact about the group. These two cylinders cannot share a
single cut at this position. This set of locks has no common blade. The
locksmith may still cut separate keys, or rearrange the cylinders, or change the
client's requirement; but within the stated problem, the answer has arrived.
The chapter's discriminating action has the same form. It does not merely lack
a certificate today. It carries a conflict that prevents the certificate it
would need. Where the telescoping action receives the cut key, the
discriminating action receives the incompatibility chart.

The final derivative reading also has a locksmith's analogue, though the shop
keeps it modest. A finished key has an exact profile. Once the cuts meet their
depths, the blade has a reading: this much lift here, that much lift there, no
leftover gap to explain by appeal to a smoother scale. The depth gauge does not
summon a continuum in order to certify the result. It reads the finite notches.
The lock confirms the same reading by turning. Exactness here is not a
philosophy added to the blade. It is the alignment of cut, pin, and shear line.
The chapter's finite Frechet derivative should be heard in that register. It is
the exact reading the apparatus earns once the finite cuts have been made.

The profile is also unique in the only sense the lock needs. If two blades open
the same cylinder by the same pin lifts, then the cylinder cannot distinguish
their working part. The bow may differ, the metal may shine differently, and
one blade may have a smoother edge, but the operative reading has already been
fixed by the pins. The lock has no extra slot in which a second profile could
hide. The chapter's uniqueness claim has the same modesty: once the finite
matching condition has fixed the operative reading, another alleged reading
must agree where the apparatus can test it.

This is why the master key belongs at the close of the chapter. It gathers the
chapter's claims without enlarging them. Exact cut, finite quotient, forced
work order, cumulative rack, compatible group, incompatible group, produced
key, tested key: these are enough. The locksmith does not need to predict a
larger theory of every possible door. He needs to show which finite pattern the
present locks can read, which pattern opens them, and which demanded pattern
cannot exist. His craft is not less rigorous because it is local. It is rigorous
because it knows the local test and submits the brass to it.

So the chapter closes at the bench. The apparatus stands beside the locksmith,
not above him. It does not assume the key it needs. It cuts the key, places it
in the cylinder, and records the turn. Where the pins conflict, it records the
conflict with the same exactness. Arrival, not approach; depth pattern, not
ornament; produced artifact, not promise; honest refusal, not postponed hope:
that is the constructive boundary in brass. A master key, when it exists, is a
made thing. When it cannot exist, the impossibility is also a made thing, read
from the finite pins and carried back to the record.
